\begin{thebibliography}{9}
\bibitem{cqCCM} Alain Connes, Caterina Consani, Henri Moscovici,
\emph{Zeta zeros and prolate wave operators: Semilocal adelic operators},
Annals of Functional Analysis (2024),
\href{https://doi.org/10.1007/s43034-024-00388-z}{doi:10.1007/s43034-024-00388-z}.
Primary author's proof PDF, 38 pages, typeset 12 September 2024;
Section 3.6, pages 17--19, Lemma 3.3, equations (3.16)--(3.22),
and Proposition 3.6.
\href{https://alainconnes.org/wp-content/uploads/Zeta-zeros-and-prolateproofs-final-2024.pdf}{Author-hosted primary proof PDF}.
The primary sources
\href{https://arxiv.org/abs/2310.18423v1}{arXiv:2310.18423v1}
and \href{https://arxiv.org/abs/2310.18423v2}{v2} were also inspected:
v1 source lines 683--704, and v2
\texttt{mainc2m24fine.tex}, lines 668--689.
\bibitem{cqCM} Alain Connes and Matilde Marcolli,
\emph{Noncommutative Geometry, Quantum Fields and Motives},
AMS Colloquium Publications 55, American Mathematical Society, 2008.
Chapter 2, Section 6.1, Proposition 2.24 and proof,
printed pages 379--381, with preceding map and boundary conditions
on printed pages 377--378, equations (2.180)--(2.199).
\bibitem{cqRT} Brad Rodgers and Terence Tao,
\emph{The de Bruijn--Newman constant is non-negative},
Forum of Mathematics, Pi 8 (2020), e6,
\href{https://doi.org/10.1017/fmp.2020.6}{doi:10.1017/fmp.2020.6},
Section 1, definition of $H_t$ and $\Phi$ and the identity
$H_0(z)=\xi(1/2+iz/2)/8$.
\end{thebibliography}
